% Blueprint content, shared by the web and pdf versions.
%
% A condensed §2 of thesis/original/main.tex. "Thesis (n)" cites equation
% (n) of the thesis as compiled (the published PDF). Remarks marked
% "Finding" mirror FINDINGS.md.
%
% \leanok on a statement: stated in Lean. \leanok in a proof: proved,
% no sorry. Keep in sync with PROGRESS.md (scripts/blueprint_check.py).

\chapter*{Introduction}

A portfolio is split into cohorts $j = 1, \dots, J$, observed over periods
$t = 1, \dots, T$. We want the true key ratio $\mu_j$ of each cohort. Two
estimates are available: the cohort's own estimate $\hmu_j$, which is
noisy, and the portfolio estimate $\hmu$, which is stable but ignores
differences between cohorts. Credibility theory blends them:
\[
  \hmu_j^b = z_j \hmu_j + (1 - z_j)\hmu .
\]
The classical Bühlmann--Straub model assumes the observations of different
cohorts are uncorrelated. The thesis drops that assumption: within a
period, the noise of cohorts $i$ and $j$ has correlation $\rho_{ij}$. The
weights inside $\hmu_j$ and $\hmu$, and the factor $z_j$, are each chosen to
minimise an error variance (\eqth{3}--(5)). The result is the credibility
factor
\[
  z_j = \frac{(1 - b_j)\tau^2}{\tau^2 + \sigma_{jj}^2 - s^2},
\]
with GLS portfolio weights $\bb = K^{-1}\bu / \bu^T K^{-1} \bu$. It reduces
to Bühlmann--Straub when the correlations vanish.

Second moments are taken as given numbers throughout: the statements are
about the formulas the moments must satisfy, not (yet) about random
variables.

% ---------------------------------------------------------------------------
\chapter{The best credibility factor}\label{ch:blend}
% Thesis §2.1.1, eqs. (6)-(7). Lean: Credibility/Blend.lean

Fix a cohort $j$ and suppose the blend is unbiased. Its error variance is a
quadratic in $z$ whose coefficients are six second moments.

\begin{definition}[Blend moments and error variance]\label{def:errVar}
  \lean{Credibility.BlendMoments, Credibility.BlendMoments.errVar}
  \leanok
  Given $\Var(\hmu_j)$, $\Var(\hmu)$, $\Var(\mu_j)$, $\cov(\hmu_j, \hmu)$,
  $\cov(\hmu_j, \mu_j)$ and $\cov(\hmu, \mu_j)$, let
  \begin{align*}
    E(z) = \Var(z\hmu_j + (1-z)\hmu - \mu_j)
      &= z^2\Var(\hmu_j) + (1-z)^2\Var(\hmu) + \Var(\mu_j)
         + 2z(1-z)\cov(\hmu_j, \hmu) \\
      &\quad - 2z\cov(\hmu_j, \mu_j) - 2(1-z)\cov(\hmu, \mu_j),
  \end{align*}
  as in \eqth{6}.
\end{definition}

\begin{definition}[Curvature and the general credibility factor]\label{def:bestZ}
  \lean{Credibility.BlendMoments.gap, Credibility.BlendMoments.bestZ}
  \leanok
  \uses{def:errVar}
  Let $g = \Var(\hmu_j - \hmu) = \Var(\hmu) - 2\cov(\hmu_j, \hmu) + \Var(\hmu_j)$
  and
  \[
    z^* = \frac{\Var(\hmu) - \cov(\hmu_j, \hmu) + \cov(\hmu_j, \mu_j) - \cov(\hmu, \mu_j)}{g},
  \]
  the general credibility factor \eqth{7}.
\end{definition}

\begin{lemma}[Vertex form]\label{lem:errVar_eq}
  \lean{Credibility.BlendMoments.errVar_eq}
  \leanok
  \uses{def:bestZ}
  If $g \neq 0$ then $E(z) = E(z^*) + g\,(z - z^*)^2$ for all $z$.
\end{lemma}
\begin{proof}
  \leanok
  Both sides are quadratics in $z$ with leading coefficient $g$; since
  $g z^*$ equals the numerator of $z^*$, their linear and constant terms
  agree too.
\end{proof}

\begin{theorem}[Eq.\ (7) is optimal]\label{thm:bestZ_optimal}
  \lean{Credibility.BlendMoments.bestZ_optimal}
  \leanok
  \uses{lem:errVar_eq}
  If $g > 0$ then $E(z^*) \le E(z)$ for all $z$.
\end{theorem}
\begin{proof}
  \leanok
  $g(z - z^*)^2 \ge 0$.
\end{proof}

\begin{theorem}[\ldots and unique]\label{thm:bestZ_unique}
  \lean{Credibility.BlendMoments.bestZ_unique}
  \leanok
  \uses{lem:errVar_eq}
  If $g > 0$ and $E(z) \le E(z^*)$ then $z = z^*$.
\end{theorem}
\begin{proof}
  \leanok
  Then $g(z - z^*)^2 \le 0$, so $(z - z^*)^2 = 0$.
\end{proof}

% ---------------------------------------------------------------------------
\chapter{The correlated data model}\label{ch:model}
% Thesis §2.3-2.4, eqs. (15)-(22). Lean: Credibility/Model.lean

\begin{definition}[Data model]\label{def:Model}
  \lean{Credibility.Model}
  \leanok
  The model \eqth{15} is
  \[
    \mu_j = \mu + \tau\varepsilon_j, \qquad
    \mu_{jt} = \mu_j + \sigma_{jt}\varepsilon_{jt},
  \]
  with $\varepsilon_j$, $\varepsilon_{jt}$ centred, of unit variance,
  $\varepsilon_j$ uncorrelated with everything else, and
  $\cov(\varepsilon_{it}, \varepsilon_{js}) = \rho_{ij}\delta_{ts}$.
  Its parameters are $\tau^2 > 0$, noise levels $\sigma_{jt} > 0$, and a
  symmetric correlation table $\rho$ with $\rho_{jj} = 1$.
\end{definition}

\begin{definition}[Observation covariance]\label{def:obsCov}
  \lean{Credibility.Model.obsCov}
  \leanok
  \uses{def:Model}
  By the law of total covariance,
  $\bar K_{ij}^{ts} = \cov(\mu_{it}, \mu_{js})
    = \sigma_{it}\sigma_{js}\rho_{ij}\delta_{ts} + \tau^2\delta_{ij}$
  \eqth{16}.
\end{definition}

The estimators are weighted averages \eqth{17},
$\hmu_j = \sum_t a_{jt}\mu_{jt}$ and $\hmu = \sum_j b_j\hmu_j$, unbiased
when the weights sum to one. The error variance of $\hmu_j$ is
$\sigma_{jj}^2$ \eqth{19}, a special case of the following.

\begin{definition}[Cohort weights and $\sigma_{ij}^2$]\label{def:sigma2}
  \lean{Credibility.Model.a, Credibility.Model.σ2}
  \leanok
  \uses{def:Model, def:invVarWeight}
  The cohort weights $a_{jt}$ are the inverse-variance weights of
  $\sigma_{j1}^2, \dots, \sigma_{jT}^2$ (Definition~\ref{def:invVarWeight},
  chosen in Chapter~\ref{ch:weights}), and
  $\sigma_{ij}^2 = \sum_t a_{it}a_{jt}\sigma_{it}\sigma_{jt}$ \eqth{18}.
\end{definition}

\begin{definition}[The covariance matrix $K$]\label{def:K}
  \lean{Credibility.Model.K}
  \leanok
  \uses{def:obsCov, def:sigma2}
  $K_{ij} = \cov(\hmu_i, \hmu_j) = \sum_{t,s} a_{it}a_{js}\bar K_{ij}^{ts}$
  \eqth{21}.
\end{definition}

\begin{lemma}[Entries of $K$]\label{lem:K_apply}
  \lean{Credibility.Model.K_apply}
  \leanok
  \uses{def:K}
  If $T \ge 1$ then $K_{ij} = \sigma_{ij}^2\rho_{ij} + \tau^2\delta_{ij}$
  (\eqth{21}, last line; also in (27)).
\end{lemma}
\begin{proof}
  \leanok
  Split the double sum. The $\delta_{ts}$ term collapses to $\sigma_{ij}^2\rho_{ij}$.
  The $\tau^2\delta_{ij}$ term is $\tau^2\delta_{ij}(\sum_t a_{it})(\sum_s a_{js})
  = \tau^2\delta_{ij}$, since inverse-variance weights sum to one
  (Lemma~\ref{lem:sum_invVarWeight}; this is where $T \ge 1$ is needed).
  \uses{lem:sum_invVarWeight}
\end{proof}

\begin{lemma}\label{lem:K_symm}
  \lean{Credibility.Model.K_symm}
  \leanok
  \uses{def:K}
  $K$ is symmetric.
\end{lemma}
\begin{proof}
  \leanok
  $\bar K_{ij}^{ts} = \bar K_{ji}^{st}$ by symmetry of $\rho$; swap the
  summation indices.
\end{proof}

\begin{lemma}[$K$ is positive definite]\label{lem:K_posDef}
  \lean{Credibility.Model.K_posDef}
  \leanok
  \uses{lem:K_apply, lem:K_symm}
  If $T \ge 1$ and the matrix $(\rho_{ij})$ is positive semidefinite, then
  $K$ is positive definite.
\end{lemma}
\begin{proof}
  \leanok
  \uses{lem:K_apply}
  For $x \in \mathbb R^J$,
  $x^T K x = \sum_t y_t^T R\, y_t + \tau^2\lVert x\rVert^2$ where
  $R = (\rho_{ij})$ and $(y_t)_j = a_{jt}\sigma_{jt}x_j$. The first sum is
  $\ge 0$ and the second is $> 0$ for $x \ne 0$.
\end{proof}

\begin{remark}[Finding 1: positive \emph{semi}definite is not enough]\label{rem:posdef}
  \uses{lem:K_posDef}
  After \eqth{20} the thesis says ``since $K$ is a covariance matrix it is
  positive definite''. A covariance matrix is only positive semidefinite.
  Here $K$ is positive definite because $\tau^2 > 0$ and the correlations
  form a positive-semidefinite matrix; the latter is an extra hypothesis of
  Lemma~\ref{lem:K_posDef}, automatic when $\rho$ really is a correlation
  matrix.
\end{remark}

The blend of Chapter~\ref{ch:blend} needs six moments. With
$\upsilon = K\bb$, so $\upsilon_j = \cov(\hmu, \hmu_j)$, the thesis computes
them for the linear estimators.

\begin{definition}[Blend moments of the model]\label{def:moments}
  \lean{Credibility.Model.moments}
  \leanok
  \uses{def:K, def:bestZ, def:Model.b}
  For cohort $j$: $\Var(\hmu_j) = \tau^2 + \sigma_{jj}^2$,
  $\Var(\hmu) = s^2$, $\Var(\mu_j) = \tau^2$,
  $\cov(\hmu_j, \hmu) = \upsilon_j$, $\cov(\hmu_j, \mu_j) = \tau^2$ and
  $\cov(\hmu, \mu_j) = b_j\tau^2$ (§2.4).
\end{definition}

\begin{theorem}[Credibility for linear estimators]\label{thm:bestZ_linear}
  \lean{Credibility.Model.bestZ_linear}
  \leanok
  \uses{def:moments}
  \[
    z_j = \frac{s^2 - \upsilon_j + \tau^2 - b_j\tau^2}{s^2 - 2\upsilon_j + \tau^2 + \sigma_{jj}^2}
    \qquad\text{\eqth{22}.}
  \]
\end{theorem}
\begin{proof}
  \leanok
  Substitute Definition~\ref{def:moments} into \eqth{7}.
\end{proof}

% ---------------------------------------------------------------------------
\chapter{The best unbiased linear weights}\label{ch:weights}
% Thesis §2.5, eqs. (23)-(27). Lean: Credibility/Weights.lean, Model.lean

\section{Cohort level: inverse-variance weights}

Minimising $\Var(\hmu_j - \mu_j) = \sum_t a_{jt}^2\sigma_{jt}^2$ subject to
$\sum_t a_{jt} = 1$ gives inverse-variance weights. The thesis uses a
Lagrange multiplier; the Lean proof completes the square instead.

\begin{definition}[Inverse-variance weights]\label{def:invVarWeight}
  \lean{Credibility.invVarWeight}
  \leanok
  $a_t = \sigma_t^{-2} / \sum_s \sigma_s^{-2}$ \eqth{23}.
\end{definition}

\begin{lemma}\label{lem:sum_inv_pos}
  \lean{Credibility.sum_inv_pos}
  \leanok
  If $T \ge 1$ and all $\sigma_t^2 > 0$ then $\sum_s \sigma_s^{-2} > 0$.
\end{lemma}
\begin{proof}
  \leanok
  A nonempty sum of positive terms.
\end{proof}

\begin{lemma}\label{lem:sum_invVarWeight}
  \lean{Credibility.sum_invVarWeight}
  \leanok
  \uses{def:invVarWeight, lem:sum_inv_pos}
  Under the same hypotheses, $\sum_t a_t = 1$.
\end{lemma}
\begin{proof}
  \leanok
  \uses{lem:sum_inv_pos}
  The numerators sum to the denominator, which is nonzero.
\end{proof}

\begin{lemma}[Minimal value]\label{lem:invVarWeight_value}
  \lean{Credibility.invVarWeight_value}
  \leanok
  \uses{def:invVarWeight, lem:sum_inv_pos}
  $\sum_t a_t^2\sigma_t^2 = \bigl(\sum_s \sigma_s^{-2}\bigr)^{-1}$.
\end{lemma}
\begin{proof}
  \leanok
  Each term is $\sigma_t^{-2} / (\sum_s \sigma_s^{-2})^2$.
\end{proof}

\begin{theorem}[Eq.\ (23) is optimal]\label{thm:invVarWeight_optimal}
  \lean{Credibility.invVarWeight_optimal}
  \leanok
  \uses{lem:invVarWeight_value}
  For any weights $c$ with $\sum_t c_t = 1$,
  $\sum_t a_t^2\sigma_t^2 \le \sum_t c_t^2\sigma_t^2$.
\end{theorem}
\begin{proof}
  \leanok
  \uses{lem:invVarWeight_value}
  $0 \le \sum_t (c_t - a_t)^2\sigma_t^2
    = \sum_t c_t^2\sigma_t^2 - 2\sum_t c_t a_t\sigma_t^2 + \sum_t a_t^2\sigma_t^2$,
  and $a_t\sigma_t^2$ is constant in $t$, so the cross term is
  $2\sum_t c_t a_t\sigma_t^2 = 2(\sum_s\sigma_s^{-2})^{-1}$, twice the
  minimal value.
\end{proof}

\section{Portfolio level: GLS weights}

The estimates $\hmu_j$ are correlated, so inverse-variance weighting does
not apply. Minimising $s^2 = \bb^T K \bb$ subject to $\bu^T\bb = 1$
($\bu$ the vector of ones) gives generalised least squares weights.

\begin{definition}[GLS weights]\label{def:glsWeight}
  \lean{Credibility.invTotal, Credibility.glsWeight}
  \leanok
  Write $K^{-1}_{**} = \bu^T K^{-1}\bu$ and
  $\bb = K^{-1}\bu / K^{-1}_{**}$, i.e.\ $b_j = K^{-1}_{j*}/K^{-1}_{**}$
  (the display before \eqth{24}, and (25)).
\end{definition}

\begin{lemma}\label{lem:invTotal_pos}
  \lean{Credibility.invTotal_pos}
  \leanok
  \uses{def:glsWeight}
  If $J \ge 1$ and $K$ is positive definite then $K^{-1}_{**} > 0$.
\end{lemma}
\begin{proof}
  \leanok
  $K^{-1}$ is positive definite and $\bu \ne 0$.
\end{proof}

\begin{lemma}\label{lem:sum_glsWeight}
  \lean{Credibility.sum_glsWeight}
  \leanok
  \uses{lem:invTotal_pos}
  Then $\sum_j b_j = 1$.
\end{lemma}
\begin{proof}
  \leanok
  \uses{lem:invTotal_pos}
  $\sum_j (K^{-1}\bu)_j = K^{-1}_{**}$, which is nonzero.
\end{proof}

\begin{lemma}\label{lem:glsWeight_eq_smul}
  \lean{Credibility.glsWeight_eq_smul}
  \leanok
  \uses{def:glsWeight}
  $\bb = (K^{-1}_{**})^{-1}\, K^{-1}\bu$ as vectors.
\end{lemma}
\begin{proof}
  \leanok
  Definition.
\end{proof}

\begin{lemma}[$\upsilon = s^2\bu$]\label{lem:mulVec_glsWeight}
  \lean{Credibility.mulVec_glsWeight}
  \leanok
  \uses{lem:glsWeight_eq_smul}
  If $K$ is positive definite then
  $\upsilon = K\bb = \bu / K^{-1}_{**}$: every cohort estimate has the same
  covariance with the portfolio estimate, namely $\Var(\hmu)$.
\end{lemma}
\begin{proof}
  \leanok
  \uses{lem:glsWeight_eq_smul}
  $K$ is invertible, and $K K^{-1}\bu = \bu$.
\end{proof}

\begin{theorem}[Minimal portfolio variance]\label{thm:glsWeight_value}
  \lean{Credibility.glsWeight_value}
  \leanok
  \uses{lem:mulVec_glsWeight, lem:sum_glsWeight}
  If $J \ge 1$ and $K$ is positive definite then
  $s^2 = \bb^T K\bb = 1/K^{-1}_{**}$ \eqth{24}.
\end{theorem}
\begin{proof}
  \leanok
  \uses{lem:mulVec_glsWeight, lem:sum_glsWeight}
  $\bb^T\upsilon = (\sum_j b_j)/K^{-1}_{**}$.
\end{proof}

\begin{theorem}[GLS weights are optimal]\label{thm:glsWeight_optimal}
  \lean{Credibility.glsWeight_optimal}
  \leanok
  \uses{thm:glsWeight_value}
  If $K$ is positive definite and $\sum_j c_j = 1$ then
  $\bb^T K\bb \le c^T K c$.
\end{theorem}
\begin{proof}
  \leanok
  \uses{thm:glsWeight_value, lem:mulVec_glsWeight}
  $0 \le (c - \bb)^T K (c - \bb) = c^TKc - 2c^TK\bb + \bb^TK\bb$, using
  that $K$ is symmetric. By Lemma~\ref{lem:mulVec_glsWeight},
  $c^TK\bb = (\sum_j c_j)s^2 = s^2 = \bb^TK\bb$.
\end{proof}

\section{The proposed credibility factor}

\begin{definition}[Portfolio weights of the model]\label{def:Model.b}
  \lean{Credibility.Model.b, Credibility.Model.s2}
  \leanok
  \uses{def:K, def:glsWeight}
  $b_j$ are the GLS weights of the model's $K$, and $s^2 = 1/K^{-1}_{**}$.
\end{definition}

\begin{lemma}[The curvature]\label{lem:gap_eq}
  \lean{Credibility.Model.gap_eq}
  \leanok
  \uses{def:moments, lem:mulVec_glsWeight}
  If $K$ is positive definite then
  $\Var(\hmu_j - \hmu) = \tau^2 + \sigma_{jj}^2 - s^2$.
\end{lemma}
\begin{proof}
  \leanok
  \uses{lem:mulVec_glsWeight}
  The denominator of \eqth{22} with $\upsilon_j = s^2$.
\end{proof}

\begin{theorem}[The proposed credibility factor]\label{thm:bestZ_eq}
  \lean{Credibility.Model.bestZ_eq}
  \leanok
  \uses{thm:bestZ_linear, lem:mulVec_glsWeight, lem:gap_eq}
  If $K$ is positive definite then
  \[
    z_j = \frac{(1-b_j)\tau^2}{\tau^2 + \sigma_{jj}^2 - s^2}
        = \frac{\cov(\hmu_j - \hmu, \mu_j)}{\Var(\hmu_j - \hmu)}
    \qquad\text{\eqth{26}.}
  \]
\end{theorem}
\begin{proof}
  \leanok
  \uses{thm:bestZ_linear, lem:mulVec_glsWeight}
  Put $\upsilon_j = s^2$ (Lemma~\ref{lem:mulVec_glsWeight}) into
  Theorem~\ref{thm:bestZ_linear}: the numerator becomes $(1-b_j)\tau^2$
  and the denominator $\tau^2 + \sigma_{jj}^2 - s^2$.
\end{proof}

Together, \eqth{23}, (18), (21), (24), (25) and (26) are the proposed model
summarised in \eqth{27}. The factor has a singularity where
$\tau^2 + \sigma_{jj}^2 = s^2$, i.e.\ where $\Var(\hmu_j - \hmu) = 0$: then
$\hmu_j = \hmu$ and every $z_j$ is optimal. (In Lean, division by zero is
zero, so the formula still typechecks there.)

% ---------------------------------------------------------------------------
\chapter{The Bühlmann--Straub special case}\label{ch:bs}
% Thesis §2.6, eqs. (28)-(30). Lean: Credibility/BuhlmannStraub.lean

\begin{definition}[Bühlmann--Straub assumptions]\label{def:bsModel}
  \lean{Credibility.bsModel}
  \leanok
  \uses{def:Model}
  $\sigma_{jt}^2 = \sigma^2/w_{jt}$ with volumes $w_{jt} > 0$, and
  $\rho_{ij} = \delta_{ij}$ \eqth{28}. Write $w_{j*} = \sum_t w_{jt}$.
\end{definition}

\begin{lemma}[Volume weights]\label{lem:bs_a}
  \lean{Credibility.bs_a}
  \leanok
  \uses{def:bsModel, def:sigma2}
  $a_{jt} = w_{jt}/w_{j*}$ \eqth{29}.
\end{lemma}
\begin{proof}
  \leanok
  $\sigma_{jt}^{-2} = w_{jt}/\sigma^2$; the factor $\sigma^{-2}$ cancels.
\end{proof}

\begin{lemma}\label{lem:bs_σ2}
  \lean{Credibility.bs_σ2}
  \leanok
  \uses{lem:bs_a}
  If $T \ge 1$, $\sigma_{jj}^2 = \sigma^2/w_{j*}$ \eqth{30}.
\end{lemma}
\begin{proof}
  \leanok
  \uses{lem:bs_a}
  $\sum_t (w_{jt}/w_{j*})^2\,\sigma^2/w_{jt} = \sigma^2 w_{j*}/w_{j*}^2$.
\end{proof}

\begin{lemma}[$K$ is diagonal]\label{lem:bs_K}
  \lean{Credibility.bs_K}
  \leanok
  \uses{lem:bs_σ2, lem:K_apply}
  If $T \ge 1$, $K = \operatorname{diag}(\tau^2 + \sigma^2/w_{j*})$.
\end{lemma}
\begin{proof}
  \leanok
  \uses{lem:bs_σ2, lem:K_apply}
  Lemma~\ref{lem:K_apply} with $\rho_{ij} = \delta_{ij}$.
\end{proof}

\begin{definition}[Classical credibility factor]\label{def:bsZ}
  \lean{Credibility.bsZ}
  \leanok
  $z_j^B = w_{j*}/(w_{j*} + \sigma^2/\tau^2)$, so that
  $K_{jj} = \tau^2/z_j^B$.
\end{definition}

\begin{lemma}\label{lem:bs_Kinv_one}
  \lean{Credibility.bs_Kinv_one}
  \leanok
  \uses{lem:bs_K, def:bsZ}
  If $T \ge 1$, $K^{-1}u = (z_j^B/\tau^2)_j$.
\end{lemma}
\begin{proof}
  \leanok
  \uses{lem:bs_K}
  $K^{-1} = \operatorname{diag}\bigl(1/(\tau^2 + \sigma^2/w_{j*})\bigr)
  = \operatorname{diag}(z_j^B/\tau^2)$.
\end{proof}

\begin{lemma}\label{lem:bs_b}
  \lean{Credibility.bs_b}
  \leanok
  \uses{lem:bs_Kinv_one, def:Model.b}
  If $T \ge 1$, $b_j = z_j^B/z_*^B$ where $z_*^B = \sum_i z_i^B$.
\end{lemma}
\begin{proof}
  \leanok
  \uses{lem:bs_Kinv_one}
  By Lemma~\ref{lem:bs_Kinv_one}, $K^{-1}_{j*} = z_j^B/\tau^2$
  and $K^{-1}_{**} = z_*^B/\tau^2$.
\end{proof}

\begin{theorem}[Reduction to Bühlmann--Straub]\label{thm:bs_z}
  \lean{Credibility.bs_z}
  \leanok
  \uses{thm:bestZ_eq, lem:bs_b, lem:bs_σ2}
  If $T \ge 1$ and $J \ge 2$, the proposed factor is $z_j = z_j^B$.
\end{theorem}
\begin{proof}
  \leanok
  \uses{thm:bestZ_eq, lem:bs_b, lem:bs_σ2, lem:bs_Kinv_one}
  $K$ is diagonal with positive entries, hence positive definite. Then
  $s^2 = \tau^2/z_*^B = (\tau^2 + \sigma_{jj}^2)\,b_j$, so
  \[
    z_j = \frac{(1-b_j)\tau^2}{(\tau^2+\sigma_{jj}^2)(1-b_j)}
        = \frac{\tau^2}{\tau^2+\sigma_{jj}^2} = z_j^B .
  \]
  Cancelling $1 - b_j$ needs $b_j < 1$, which holds when $J \ge 2$
  because every $z_i^B > 0$.
\end{proof}

\begin{remark}[Finding 2: at least two cohorts]\label{rem:two_cohorts}
  \uses{thm:bs_z}
  With $J = 1$ we get $b = 1$ and $s^2 = \tau^2 + \sigma^2$, so
  \eqth{26} is $0/0$. The thesis never assumes $J \ge 2$, but the
  cancellation in the last step of §2.6 needs it.
\end{remark}

% ---------------------------------------------------------------------------
\chapter{Two cohorts with time-constant weights}\label{ch:two}
% Thesis §2.7, eqs. (31)-(34). Lean: Credibility/TwoCohorts.lean

With $J = 2$, $w_{jt} = w_j$ and $\sigma_{jt}^2 = \sigma^2/w_j$, we get
$a_{jt} = 1/T$ and $\sigma_{ij}^2 = \sigma_i\sigma_j$ with
$\sigma_i = \sigma/(T w_i)$. Write $\rho = \rho_{12}$ and
$\kappa_i = \sigma_i^2/\tau^2$. The Lean statements avoid square roots:
$\sqrt{\kappa_1\kappa_2}\,\rho = \sigma_1\sigma_2\rho/\tau^2$.

\begin{definition}[Two-cohort quantities]\label{def:two}
  \lean{Credibility.TwoCohorts.K, Credibility.TwoCohorts.D,
        Credibility.TwoCohorts.s2, Credibility.TwoCohorts.b1}
  \leanok
  \[
    K = \begin{pmatrix} \sigma_1^2 + \tau^2 & \sigma_1\sigma_2\rho \\
                        \sigma_1\sigma_2\rho & \sigma_2^2 + \tau^2 \end{pmatrix},
    \qquad D = \sigma_1^2 + \sigma_2^2 + 2\tau^2 - 2\sigma_1\sigma_2\rho,
  \]
  $s^2 = \bigl((\sigma_1^2+\tau^2)(\sigma_2^2+\tau^2) - \sigma_1^2\sigma_2^2\rho^2\bigr)/D$
  and $b_1 = (\sigma_2^2 + \tau^2 - \sigma_1\sigma_2\rho)/D$.
\end{definition}

\begin{lemma}\label{lem:D_pos}
  \lean{Credibility.TwoCohorts.D_pos}
  \leanok
  \uses{def:two}
  If $\sigma_1, \sigma_2, \tau^2 > 0$ and $\rho \le 1$ then $D > 0$.
\end{lemma}
\begin{proof}
  \leanok
  $D = (\sigma_1 - \sigma_2)^2 + 2\tau^2 + 2\sigma_1\sigma_2(1 - \rho)$.
\end{proof}

\begin{theorem}[GLS weights for two cohorts]\label{thm:gls_two}
  \lean{Credibility.TwoCohorts.gls_two}
  \leanok
  \uses{lem:D_pos, def:glsWeight}
  If moreover $|\rho| < 1$, the GLS weight of cohort 1 is $b_1$ and
  $1/K^{-1}_{**} = s^2$ (\eqth{31} and the display before it).
\end{theorem}
\begin{proof}
  \leanok
  \uses{lem:D_pos}
  $\det K = (\sigma_1^2+\tau^2)(\sigma_2^2+\tau^2) - \sigma_1^2\sigma_2^2\rho^2 > 0$
  since $\rho^2 < 1$. Invert $K$ by the adjugate formula; the entries of
  $K^{-1}\bu$ are $(\sigma_2^2+\tau^2-\sigma_1\sigma_2\rho)/\det K$ and
  $(\sigma_1^2+\tau^2-\sigma_1\sigma_2\rho)/\det K$, which sum to
  $D/\det K$.
\end{proof}

\begin{theorem}[Two-cohort credibility factor]\label{thm:z1_eq}
  \lean{Credibility.TwoCohorts.z1_eq}
  \leanok
  \uses{thm:gls_two}
  Away from the singularity $\tau^2 + \sigma_1^2 = \sigma_1\sigma_2\rho$,
  \[
    z_1 = \frac{(1 - b_1)\tau^2}{\tau^2 + \sigma_1^2 - s^2}
        = \frac{\tau^2}{\tau^2 + \sigma_1^2 - \sigma_1\sigma_2\rho}
        = \frac{1}{1 + \kappa_1 - \sqrt{\kappa_1\kappa_2}\,\rho}
    \qquad\text{\eqth{32}.}
  \]
\end{theorem}
\begin{proof}
  \leanok
  \uses{lem:D_pos}
  With $A = \sigma_1^2 + \tau^2$ and $c = \sigma_1\sigma_2\rho$:
  $1 - b_1 = (A - c)/D$ and $A D - \det K = (A - c)^2$, so
  $\tau^2 + \sigma_1^2 - s^2 = (A - c)^2/D$. Cancel $(A - c)/D$.
\end{proof}

\begin{theorem}[Weights from credibility factors]\label{thm:b1_eq_z}
  \lean{Credibility.TwoCohorts.b1_eq_z}
  \leanok
  \uses{thm:gls_two, thm:z1_eq}
  Away from both singularities, with $z_i = \tau^2/(\tau^2 + \sigma_i^2 - \sigma_1\sigma_2\rho)$,
  \[
    b_1 = \frac{\kappa_2 + 1 - \sqrt{\kappa_1\kappa_2}\,\rho}
               {\kappa_1 + \kappa_2 + 2 - 2\sqrt{\kappa_1\kappa_2}\,\rho}
        = \frac{z_1}{z_1 + z_2}.
  \]
\end{theorem}
\begin{proof}
  \leanok
  \uses{lem:D_pos}
  With $A_i = \tau^2 + \sigma_i^2 - \sigma_1\sigma_2\rho$ we have $z_i = \tau^2/A_i$,
  so $z_1/(z_1+z_2) = A_2/(A_1 + A_2)$, and $A_1 + A_2 = D$.
\end{proof}

\begin{remark}[Finding 4: only for two cohorts]\label{rem:b1_eq_z}
  \uses{thm:b1_eq_z}
  The second equality is not in the published thesis; it comes from the
  defence slides. It is the Bühlmann--Straub form $b_j = z_j/z_*$, which
  therefore survives correlation for any binary split of a portfolio
  (e.g.\ two segments of a rating factor). For three or more correlated
  cohorts it fails in general (numerically, e.g.\ $J = 3$ with random
  $\sigma_j$ and $\rho_{ij}$), even with equal correlations
  $\rho_{ij} = \rho$. What carries over is equal \emph{covariances}
  (Theorem~\ref{thm:b_eq_of_equalCov}).
\end{remark}

\paragraph{More than two cohorts: equal covariances.}
With two cohorts there is a single covariance $\sigma_1\sigma_2\rho$. The
right generalisation is that all cohort pairs share one covariance
$K_{ij} = \sigma_{ij}^2\rho_{ij} = c$ ($i \ne j$). Then
$K = \operatorname{diag}(d_j) + c\,uu^\top$ with
$d_j = \tau^2 + \sigma_{jj}^2 - c$, and the rank-one part lies along $u$.

\begin{lemma}\label{lem:Kinv_one_of_equalCov}
  \lean{Credibility.Model.Kinv_one_of_equalCov}
  \leanok
  \uses{lem:K_apply, def:K}
  If $K$ is positive definite, $K_{ij} = c$ for $i \ne j$ and all $d_j \ne 0$,
  then with $v_j = 1/d_j$: $1 + c\,v_* \ne 0$ and
  $K^{-1}u = v/(1 + c\,v_*)$.
\end{lemma}
\begin{proof}
  \leanok
  $(Kv)_i = d_i v_i + c\,v_* = 1 + c\,v_*$, so $Kv = (1 + c\,v_*)u$. Apply
  $K^{-1}$; if $1 + c\,v_* = 0$ then $v = 0$, impossible.
\end{proof}

\begin{theorem}[Weights from credibility factors, any $J$]\label{thm:b_eq_of_equalCov}
  \lean{Credibility.Model.b_eq_of_equalCov}
  \leanok
  \uses{lem:Kinv_one_of_equalCov, def:Model.b}
  Under the same hypotheses, with $z_j = \tau^2/(\tau^2 + \sigma_{jj}^2 - c)$,
  $b_j = z_j/z_*$.
\end{theorem}
\begin{proof}
  \leanok
  $b = K^{-1}u/u^\top K^{-1}u = v/v_*$ and $z = \tau^2 v$.
\end{proof}

\begin{theorem}\label{thm:bestZ_of_equalCov}
  \lean{Credibility.Model.bestZ_of_equalCov}
  \leanok
  \uses{lem:Kinv_one_of_equalCov, thm:bestZ_eq}
  If moreover $b_j \ne 1$, the proposed factor is
  $z_j = \tau^2/(\tau^2 + \sigma_{jj}^2 - c)$: eq.~\eqth{32} with
  $\sigma_1\sigma_2\rho$ replaced by $c$.
\end{theorem}
\begin{proof}
  \leanok
  \uses{thm:bestZ_eq}
  $s^2 = (1 + c\,v_*)/v_* = c + 1/v_*$, so
  $\tau^2 + \sigma_{jj}^2 - s^2 = d_j - 1/v_* = (v_* - v_j)/(v_j v_*)$ and
  $1 - b_j = (v_* - v_j)/v_*$. Cancel $v_* - v_j \ne 0$.
\end{proof}

\begin{remark}[Finding 5: equal covariances is the right condition]\label{rem:equalCov}
  \uses{thm:b_eq_of_equalCov}
  For $J = 3$ the condition is also necessary (when all $b_j \ne 0$).
  Writing $(Kb)_j = s^2$ shows that $b_j = z_j/z_*$ for all $j$ iff there is a
  $\lambda$ with $\sum_{i \ne j} (K_{ij} + \lambda)\, b_i = 0$ for every $j$.
  With $p, q, r$ the three values $K_{ij} + \lambda$ this is
  $p b_2 + q b_3 = p b_1 + r b_3 = q b_1 + r b_2 = 0$, which forces
  $p = q = r = 0$. For $J \ge 4$ it is not necessary: with equal $\sigma_j$, any
  correlation matrix with equal row sums gives $b = z/z_* = u/J$, whether or not
  the covariances are equal. Bühlmann--Straub is the case $c = 0$; two cohorts
  are the case $c = \sigma_1\sigma_2\rho$. Not yet in Lean.
\end{remark}

\begin{remark}[Finding 6: equal credibility is not grouping]\label{rem:equalZ_grouping}
  \uses{thm:b_eq_of_equalCov}
  Under equal covariances, cohorts with the same $z_j$ get the same weight
  $b_j$, so $\hmu = \sum_j b_j \hmu_j$ depends on them only through
  $\sum \hmu_j$. Pooling them loses nothing for $\hmu$. This is not risk
  grouping, though. Equal $z$ means equal reliability, not equal level:
  premiums $z\hmu_j + (1-z)\hmu$ still differ. Merging $i$ and $j$ into one
  risk group assumes $\mu_i = \mu_j$, which changes their covariance from
  $0$ to $\tau^2$ and hence $K$, every $z$ and $\hmu$.

  When merging pays off (on paper, \texttt{notes/merging-cohorts.md}): in a
  two-level model, let $d = \mathbb{E}(\mu_1-\mu_2)^2$ and
  $V = \operatorname{Var}(\hmu_1 - \hmu_2) = \sigma_1^2+\sigma_2^2-2\sigma_1\sigma_2\rho$.
  For two cohorts on their own, merging lowers the total mean squared error
  iff $d < V\tau^2/(V+\tau^2)$. Correlation enters only through $V$, and
  $\rho > 0$ makes merging harder, since the gap $\hmu_1 - \hmu_2$ is then
  measured more precisely. Equal $z$ plays no role. Under the thesis model,
  $d = 2\tau^2$, merging never pays. See Chapter~\ref{ch:merging}.
\end{remark}

\begin{corollary}[Bühlmann--Straub at $\rho = 0$]\label{cor:z1_rho_zero}
  \lean{Credibility.TwoCohorts.z1_rho_zero}
  \leanok
  \uses{thm:z1_eq}
  At $\rho = 0$, $z_1 = \tau^2/(\tau^2 + \sigma_1^2) = 1/(1 + \kappa_1)$.
\end{corollary}
\begin{proof}
  \leanok
  Immediate.
\end{proof}

Unlike Bühlmann--Straub, $z_1$ now depends on $\kappa_2$ and $\rho$ too. The
singularity sits at
$\rho^* = (\tau^2 + \sigma_1^2)/(\sigma_1\sigma_2) = (1+\kappa_1)/\sqrt{\kappa_1\kappa_2}$.

\begin{theorem}[Sign flip past the singularity]\label{thm:z1_neg_iff}
  \lean{Credibility.TwoCohorts.z1_neg_iff}
  \leanok
  \uses{thm:z1_eq}
  If $\sigma_1, \sigma_2, \tau^2 > 0$ then $z_1 < 0$ iff $\rho > \rho^*$
  (displays after \eqth{32}).
\end{theorem}
\begin{proof}
  \leanok
  $\tau^2 > 0$, so $z_1$ has the sign of its denominator
  $\sigma_1\sigma_2(\rho^* - \rho)$.
\end{proof}

\begin{theorem}[When the singularity is reachable]\label{thm:singularity_reachable_iff}
  \lean{Credibility.TwoCohorts.singularity_reachable_iff}
  \leanok
  \uses{def:two}
  If $\sigma_1, \sigma_2, \tau^2 > 0$ then $\rho^* \le 1$ iff
  $\kappa_2 \ge (1 + \kappa_1)^2/\kappa_1$.
\end{theorem}
\begin{proof}
  \leanok
  Both say $(\tau^2 + \sigma_1^2)^2 \le \sigma_1^2\sigma_2^2$, i.e.\
  $\tau^2 + \sigma_1^2 \le \sigma_1\sigma_2$, as both sides are positive.
\end{proof}

The thesis shows that $s^2/\tau^2$ \eqth{33} is maximal at $\rho = \rho^*$,
so $\hmu$ is least accurate there, and proposes clipping
$\tilde z_i = \max(\min(z_i, 1), 0)$ \eqth{34}. It notes that $z_1 > 1$
can occur, but not when. The following answers that.

\begin{theorem}[Credibility above one]\label{thm:z1_gt_one_iff}
  \lean{Credibility.TwoCohorts.z1_gt_one_iff}
  \leanok
  \uses{thm:z1_eq}
  If $\sigma_1, \sigma_2 > 0$ and $\rho < \rho^*$ then
  $z_1 > 1$ iff $\sigma_1 < \sigma_2\rho$.
\end{theorem}
\begin{proof}
  \leanok
  Below the singularity the denominator $\tau^2 + \sigma_1^2 - \sigma_1\sigma_2\rho$
  is positive, so $z_1 > 1$ iff $\sigma_1^2 < \sigma_1\sigma_2\rho$, i.e.\
  $\sigma_1 < \sigma_2\rho$.
\end{proof}

\begin{remark}[Finding 3: when $z_1 > 1$]\label{rem:gt_one}
  \uses{thm:z1_gt_one_iff}
  This sharpens §2.7. In terms of $\kappa$: $z_1 > 1$ iff
  $\sqrt{\kappa_1} < \sqrt{\kappa_2}\,\rho$: cohort 1 is the less noisy one
  and the correlation is strong enough (``noise cancellation''), so the
  optimal blend overshoots $\hmu_1$. The condition
  does not involve $\tau^2$: the Lean statement carries the hypothesis
  $\tau^2 > 0$, but the proof never uses it, so the condition is
  independent of the between-cohort variance.
\end{remark}

% ---------------------------------------------------------------------------
\chapter{The balance property}\label{ch:balance}
% Thesis §2.9. Lean: Credibility/BuhlmannStraub.lean

Blended estimates are \emph{balanced} if
$\sum_{j,t} w_{jt}\mu_{jt} = \sum_{j,t} w_{jt}\hmu_j^b$: had the insurer
charged $\hmu_j^b$ historically, the net result would be zero. With
volume weights $a_{jt} = w_{jt}/w_{j*}$ this reads
$\sum_j w_{j*}(1 - z_j)(\hmu_j - \hmu) = 0$. It fails for the proposed model
in general (the simulation study measures by how much), but holds for
Bühlmann--Straub.

\begin{theorem}[Bühlmann--Straub is balanced]\label{thm:bs_balance}
  \lean{Credibility.bs_balance}
  \leanok
  \uses{def:bsZ}
  Let $W_j > 0$, $\kappa > 0$, $z_j = W_j/(W_j + \kappa)$ and
  $\hmu = \sum_j z_j\hmu_j / \sum_j z_j$, with $J \ge 1$. Then for
  \emph{any} numbers $\hmu_j$,
  $\sum_j W_j(1 - z_j)(\hmu_j - \hmu) = 0$.
\end{theorem}
\begin{proof}
  \leanok
  $W_j(1 - z_j) = W_j\kappa/(W_j + \kappa) = \kappa z_j$, so the sum is
  $\kappa\bigl(\sum_j z_j\hmu_j - z_*\hmu\bigr) = 0$ by the definition of
  $\hmu$.
\end{proof}

This holds whether or not the model assumptions \eqth{28} are met: it is
an identity in the data.

% ---------------------------------------------------------------------------
\chapter{Extension: when merging two cohorts pays off}\label{ch:merging}
% Not in the thesis: Result A of notes/merging-cohorts.md (Case A).
% Lean: Credibility/Merging.lean

Not in the thesis; it answers the question of
Remark~\ref{rem:equalZ_grouping}. Take two cohorts on their own, as in
Chapter~\ref{ch:two}, with $\mu$ estimated from the pair. Let
$v_j = \sigma_j^2$ be the error variances, $c$ the error covariance
($c = \sigma_1\sigma_2\rho$ in Chapter~\ref{ch:two}),
$V = \operatorname{Var}(\hmu_1 - \hmu_2)$ and $d = \mathbb{E}(\mu_1-\mu_2)^2$
the true gap between the levels. An unbiased linear predictor of $\mu_1$ is
$\hmu_1 - \beta(\hmu_1 - \hmu_2)$. A working model with gap variance $x$ uses
$\beta(x) = (v_1 - c)/(V + x)$: keeping the cohorts separate (the thesis model)
is $x = 2\tau^2$, merging them ($\mu_1 = \mu_2$) is $x = 0$. Under the thesis
model $\beta(2\tau^2)$ is the weight $(1 - z_1)b_2$ of \eqth{32}.

\begin{definition}[Merging quantities]\label{def:merging}
  \lean{Credibility.Merging.V, Credibility.Merging.mse,
        Credibility.Merging.weight, Credibility.Merging.Q,
        Credibility.Merging.excessLoss}
  \leanok
  $V = v_1 + v_2 - 2c$, $Q = (v_1 - c)^2 + (v_2 - c)^2$, and the mean squared
  error of the weight $\beta$ for cohort $j$,
  \[
    \operatorname{mse}_j(\beta) = v_j - 2\beta(v_j - c) + \beta^2(V + d).
  \]
  The excess loss of the working value $d_w$ is
  $\Delta L(d_w) = \sum_{j=1}^2 \bigl(\operatorname{mse}_j(\beta_j(d_w)) -
  \operatorname{mse}_j(\beta_j(d))\bigr)$ with $\beta_j(x) = (v_j - c)/(V + x)$.
\end{definition}

\begin{lemma}[Excess-loss identity]\label{lem:mse_eq}
  \lean{Credibility.Merging.mse_eq}
  \leanok
  \uses{def:merging}
  If $V + d \ne 0$, then
  $\operatorname{mse}_j(\beta) = \operatorname{mse}_j(\beta^*) + (V + d)(\beta - \beta^*)^2$
  with $\beta^* = (v_j - c)/(V + d)$.
\end{lemma}
\begin{proof}
  Complete the square.
\end{proof}

\begin{lemma}[Closed form of the excess loss]\label{lem:excessLoss_eq}
  \lean{Credibility.Merging.excessLoss_eq}
  \leanok
  \uses{lem:mse_eq}
  If $V + d \ne 0$ and $V + d_w \ne 0$, then
  \[
    \Delta L(d_w) = Q\,(V + d)\Bigl(\frac{1}{V + d_w} - \frac{1}{V + d}\Bigr)^2.
  \]
\end{lemma}
\begin{proof}
  \uses{lem:mse_eq}
  By Lemma~\ref{lem:mse_eq} cohort $j$ contributes $(V + d)(\beta_j(d_w) - \beta_j(d))^2$,
  and $\beta_j(d_w) - \beta_j(d) = (v_j - c)\bigl(1/(V + d_w) - 1/(V + d)\bigr)$.
\end{proof}

\begin{theorem}[Result A: midpoint rule]\label{thm:merge_iff}
  \lean{Credibility.Merging.merge_iff}
  \leanok
  \uses{lem:excessLoss_eq}
  Let $\tau^2 > 0$, $V > 0$, $Q > 0$ and $d \ge 0$. Merging beats keeping the
  cohorts separate, $\Delta L(0) < \Delta L(2\tau^2)$, if and only if
  \[
    d\,(V + \tau^2) < V\tau^2, \qquad\text{i.e.}\qquad d < \frac{V\tau^2}{V + \tau^2}.
  \]
\end{theorem}
\begin{proof}
  \uses{lem:excessLoss_eq}
  The common factor $Q(V + d) > 0$ cancels. With $t = 1/(V + d)$,
  \[
    \Bigl(\frac{1}{V + 2\tau^2} - t\Bigr)^2 - \Bigl(\frac1V - t\Bigr)^2
    = \frac{4\tau^2\bigl(V\tau^2 - d(V + \tau^2)\bigr)}{V^2(V + 2\tau^2)^2(V + d)}.
  \]
  So merging wins iff $t$ lies above the midpoint of $1/V$ and $1/(V + 2\tau^2)$.
\end{proof}

\begin{remark}\label{rem:merge_reading}
  \uses{thm:merge_iff}
  The condition involves only $V$, $\tau^2$ and $d$, not $z_1$ against $z_2$.
  Correlation enters through $V = \sigma_1^2 + \sigma_2^2 - 2\sigma_1\sigma_2\rho$:
  $\rho > 0$ makes merging harder. Under the thesis model $d = 2\tau^2$, and
  $2\tau^2(V + \tau^2) > V\tau^2$, so merging never pays. Because $Q$ is a
  common factor, each cohort on its own prefers the same choice.
\end{remark}
